\documentclass[11pt]{article}
\usepackage{amsmath,amssymb,amsthm,mathtools}
\usepackage[margin=1in]{geometry}

\newtheorem{theorem}{Theorem}
\newtheorem{lemma}[theorem]{Lemma}
\newtheorem{proposition}[theorem]{Proposition}
\newtheorem{corollary}[theorem]{Corollary}
\theoremstyle{definition}
\newtheorem{definition}[theorem]{Definition}
\newtheorem{remark}[theorem]{Remark}

\DeclareMathOperator{\tr}{tr}
\DeclareMathOperator{\ord}{ord}
\DeclareMathOperator{\rank}{rank}
\DeclareMathOperator{\im}{im}
\DeclareMathOperator{\Res}{Res}
\DeclareMathOperator{\Hom}{Hom}
\newcommand{\Pq}{P_{q}}
\newcommand{\Pm}{P_{-1}}
\newcommand{\la}{\lambda}
\newcommand{\Vl}{V^{\la}}
\newcommand{\Ep}[1]{E^{+}_{#1}}
\newcommand{\Em}[1]{E^{-}_{#1}}

\title{A forward collapse for the off--hook staircase monodromy:\\
the backbone is born at $T_{|\hat\la|+1}$\\
\large (the family $\la=(2,2,1^m)$, uniformly in $m$)}
\author{Clio}
\date{2026--05--29 (prove session)}

\begin{document}
\maketitle

\begin{abstract}
For the degenerate Baxterised staircase monodromy on the Hoefsmit seminormal
module $\Vl$ of $H_n(q)$, the per--subset operators
$M(S)=\prod_k A_k(S)$ ($A_k=\Pq^{(i_k)}$ or $\Pm^{(i_k)}$) control the order law
$\ord_{x=q^2}Z_\la=\tfrac{\tau(\tau+1)}2$ through the strong vanishing statement
$M(S)\equiv0$ on $\Vl$ for $|S|<D:=\tfrac{\tau(\tau+1)}2$. Working with the
\emph{forward} (low--to--high) reading of the staircase word, we prove for the entire
off--hook family $\la=(2,2,1^m)$ ($m\ge2$, so $\tau=m-1\ge1$) an exact, $m$--independent
description of the all--$\Pq$ ``backbone''. Our main rigorous result is that the prefix
product through staircase block $B_4$ has image equal to a \emph{single} seminormal basis
line $\langle v_{T_0}\rangle$, where $T_0$ has rows $(1,2),(3,4)$ and column
$(5,6,\dots,n)$; this line lies in $\bigcap_{i=5}^{n-1}\Em i$, the common
$(-1)$--eigenspace of the tail generators. The proof is uniform in $m$: branching forces
the only $\Pq$--surviving constituent of $\Vl\!\downarrow\! S_4$ to be $V^{(2,2)}$ (because
$\la_1=2$ and the degenerate monodromy vanishes on $\tau\ge1$ constituents), and the
$q$--eigenspace of $T_1$ cuts that constituent down to the single line $\langle
v_{T_0}\rangle$. The collapse is therefore born at $T_{|\hat\la|+1}=T_5$, whose first
occurrence is the head of block $B_6$ at the fixed word--position $\binom52=10$; this
rigorously explains the empirically observed $m$--independence of the collapse location.
We then derive the no--head--insertion reduction (operator vanishing of $M(S)$ becomes a
single--vector reach for $v_{T_0}$), record the forward annihilation of $v_{T_0}$ by $T_5$
(the surviving step through $B_5$ verified exactly for $m\le5$; the conclusion
$M(\varnothing)\equiv0$ is in any case secured unconditionally by the backward Pillar~1
argument), and isolate the remaining head--insertion residual precisely (the
operator--level analogue of the merged--junction gap).
\end{abstract}

\section{Setup and conventions}

We use the conventions of \texttt{PROVE.md} (2026--05--29). $H_n(q)$ has generators
$T_1,\dots,T_{n-1}$ with $(T_i+1)(T_i-q)=0$; on $\Vl$ the complementary idempotents
\[
\Pq^{(i)}=\frac{T_i+1}{q+1},\qquad \Pm^{(i)}=\frac{q-T_i}{q+1}
\]
project onto the $q$--eigenspace $\Ep i=\ker(T_i-q)$ and the $(-1)$--eigenspace
$\Em i=\ker(T_i+1)$ respectively, so that $\im\Pq^{(i)}=\Ep i$, $\ker\Pq^{(i)}=\Em i$,
$\im\Pm^{(i)}=\Em i$, $\ker\Pm^{(i)}=\Ep i$.

\paragraph{The staircase word.} Let $N=\binom n2$ and let the staircase reduced word be
$w=B_2B_3\cdots B_n$, where block $B_k=(k-1,k-2,\dots,1)$ has length $k-1$. We index the
letters $i_1,i_2,\dots,i_N$ in this left--to--right order; the $0$--indexed position of the
head of block $B_k$ is $\binom{k-1}{2}=\sum_{j=2}^{k-1}(j-1)$. For a subset
$S\subseteq\{0,\dots,N-1\}$ of word--positions put
\[
A_k(S)=\Pm^{(i_k)}\ (k\!-\!1\in S),\qquad A_k(S)=\Pq^{(i_k)}\ (k\!-\!1\notin S),
\qquad
M(S)=A_N(S)\,A_{N-1}(S)\cdots A_1(S),
\]
so that $A_1$ (position $0$) is applied first. The image of $M(S)$ is contained in the
image of its leftmost factor $A_N$. With $S=\varnothing$ this is the degenerate monodromy:
$M(\varnothing)=(q+1)^{-N}\,\Omega$, $\Omega=\prod_k(T_{i_k}+1)$.

\paragraph{The family.} For $\la=(2,2,1^m)$ one has $\hat\la=(2,2)$, $\ell=2$,
$|\hat\la|=4$, $n=m+4$, and $\tau=\tau(\la)=m-1$, $D=\binom m2$. We freely use the
combinatorial fact (verified exactly for $m=2,3,4,5$ at $q=5/7$, and structurally
$m$--uniform) that the empirical strong vanishing $M(S)\equiv0$ on $\Vl$ holds for all
$|S|<D$; the present paper proves the structural backbone underlying it.

\paragraph{The image chain.} Set $V_0=\Vl$ and $V_j=A_j(S)\,V_{j-1}=\im\!\big(A_j\cdots
A_1\big)$. Then $M(S)\equiv0$ on $\Vl$ iff $V_N=0$, i.e.\ iff some prefix vanishes:
$V_{j-1}\subseteq\ker A_j$. Since $\ker\Pq^{(i)}=\Em i$ and $\ker\Pm^{(i)}=\Ep i$, a
prefix kill happens at step $j$ exactly when
\[
V_{j-1}\subseteq\Em{i_j}\ \ (j\!-\!1\notin S),\qquad\text{or}\qquad
V_{j-1}\subseteq\Ep{i_j}\ \ (j\!-\!1\in S).
\]

\section{The backbone lemma}

The key objects are the prefix products through blocks $B_4$ and $B_5$, i.e.\ the
$S_4=\langle T_1,T_2,T_3\rangle$ and $S_5=\langle T_1,T_2,T_3,T_4\rangle$ staircase
monodromies. Write
\[
\Pi_{S_k}=\prod_{\text{letters of }B_2\cdots B_k}\Pq^{(i)}\qquad(\text{all }\Pq),
\]
the all--$\Pq$ prefix through block $B_k$; it lies in the parabolic group algebra of
$S_k=\langle T_1,\dots,T_{k-1}\rangle$.

\begin{lemma}[Constituent vanishing]\label{lem:const}
Let $\mu\vdash k$ and let $\Pi_{S_k}$ act on the irreducible seminormal module $V^\mu$ of
$H_k(q)$. If $\tau(\mu):=\max\!\big(0,\,2\ell(\mu)-k-1\big)\ge1$, then
$\Pi_{S_k}\big|_{V^\mu}=0$.
\end{lemma}

\begin{proof}
This is the operator--level vanishing of the degenerate staircase monodromy on a
constituent with $\tau\ge1$. For multi--row $\hat\mu$ it is the proved Pillar~1 corollary
($\Omega\equiv0$ on $V^\mu$ when $\tau\ge1$). For the four constituents that occur below
($\mu\in\{(2,1,1),(1^4),(2,1,1,1),(1^5)\}$) it is in addition verified directly by exact
rational computation in the corresponding modules (dimensions $3,1,4,1$); each
$\Pi_{S_k}\big|_{V^\mu}$ is the zero matrix. We will only use these four instances.
\end{proof}

\begin{lemma}[Constituent non--vanishing]\label{lem:nonvan}
$\Pi_{S_4}\big|_{V^{(2,2)}}\neq0$ and $\Pi_{S_5}\big|_{V^{(2,2,1)}}\neq0$; in fact each has
rank $1$.
\end{lemma}

\begin{proof}
Direct exact computation in the two modules (dimensions $2$ and $5$). Both shapes have
$\tau=0$, so the trace of the degenerate monodromy is nonzero, forcing rank $\ge1$; the
ranks are computed to be exactly $1$.
\end{proof}

\begin{proposition}[Branching support]\label{prop:branch}
Restrict $\Vl\!\downarrow\! S_k$ for $k\in\{4,5\}$. Because $\la_1=2$, only partitions
$\mu\vdash k$ with at most two columns occur, i.e.\ $\mu_1\le2$. Among these,
$\Pi_{S_k}$ acts nonzero only on the constituent $\mu=(2,2,1^{k-4})$, and the multiplicity
of that constituent is
\[
c_\mu=\dim\Hom_{S_k}\!\big(V^\mu,\Vl\!\downarrow\! S_k\big)=f^{\la/\mu}
=f^{(1^{\,m-(k-4)})}=1 .
\]
\end{proposition}

\begin{proof}
The partitions of $k$ with $\mu_1\le2$ are: for $k=4$, $\{(2,2),(2,1,1),(1^4)\}$; for
$k=5$, $\{(2,2,1),(2,1,1,1),(1^5)\}$. Their $\tau$--values are $\{0,1,3\}$ and
$\{0,2,4\}$ respectively, so the unique $\tau=0$ shape is $(2,2,1^{k-4})$; by
Lemma~\ref{lem:const} all other constituents are annihilated by $\Pi_{S_k}$. The
restriction multiplicity equals the number of paths of length $n-k$ from $\mu$ to $\la$ in
Young's lattice, i.e.\ the number $f^{\la/\mu}$ of standard tableaux of the skew shape
$\la/\mu$. For $\mu=(2,2,1^{k-4})\subseteq\la=(2,2,1^m)$ the skew shape is the single
column $(1^{\,m-(k-4)})$ (the bottom cells of column $0$), which admits exactly one
increasing filling, so $c_\mu=1$.
\end{proof}

Since $\Pi_{S_k}\in\mathbb{F}[S_k]$ acts block--diagonally on the isotypic decomposition
of $\Vl\!\downarrow\! S_k$, acting as $\Pi_{S_k}\big|_{V^\mu}\otimes\mathrm{id}$ on the
$V^\mu$--isotypic component, Proposition~\ref{prop:branch} gives at once that $\im\Pi_{S_k}$
is contained in the $V^{(2,2,1^{k-4})}$--isotypic component of $\Vl\!\downarrow\! S_k$. For
$k=4$ this pins the image to a single line, with no further computation.

\begin{proposition}[The backbone line]\label{prop:line}
Let $T_0$ be the standard tableau of $\la$ with rows $(1,2),(3,4)$ and column
$(5,6,\dots,n)$. Then, uniformly in $m\ge2$,
\[
\im\Pi_{S_4}=\langle v_{T_0}\rangle .
\]
\end{proposition}

\begin{proof}
By the remark above, $\im\Pi_{S_4}$ lies in the $V^{(2,2)}$--isotypic component
$\mathcal I$ of $\Vl\!\downarrow\! S_4$. The leftmost factor of the prefix product
$\Pi_{S_4}=\Pi_{B_2B_3B_4}$ is $\Pq^{(1)}$ (the last letter of $B_4$ is $T_1$), so
$\im\Pi_{S_4}\subseteq\Ep1$. Now $\mathcal I$ is spanned by the seminormal vectors $v_T$
for which the entries $\{1,2,3,4\}$ fill the cells $(0,0),(0,1),(1,0),(1,1)$; there are
exactly two such standard tableaux of $\la$ (with multiplicity $c_{(2,2)}=f^{\la/(2,2)}
=f^{(1^m)}=1$), namely $T_0$ and the tableau with $2$ in cell $(1,0)$. Since $T_i$ is
diagonal in the seminormal basis with eigenvalue $q$ exactly when $i,i+1$ share a row, the
intersection $\Ep1\cap\mathcal I$ is spanned by those $v_T\in\mathcal I$ with $1,2$ in the
same row --- only $T_0$ qualifies. Hence $\Ep1\cap\mathcal I=\langle v_{T_0}\rangle$ is
one--dimensional, so $\im\Pi_{S_4}\subseteq\langle v_{T_0}\rangle$. By
Lemma~\ref{lem:nonvan} and Proposition~\ref{prop:branch} the rank of $\Pi_{S_4}$ on $\Vl$
is $c_{(2,2)}\cdot\rank\Pi_{S_4}|_{V^{(2,2)}}=1$, so the image is exactly $\langle
v_{T_0}\rangle$.
\end{proof}

\begin{lemma}[The line lies in $\Em5$]\label{lem:inE5}
$v_{T_0}\in\Em5$; more precisely $v_{T_0}\in\Ep1\cap\Ep3\cap\bigcap_{i=5}^{n-1}\Em i$.
\end{lemma}

\begin{proof}
Seminormal/Hoefsmit eigenvalues are read off the relative positions of $i,i+1$ in $T_0$:
if $i,i+1$ lie in the same row then $T_iv_{T_0}=q\,v_{T_0}$ (so $v_{T_0}\in\Ep i$), and if
they lie in the same column then $T_iv_{T_0}=-v_{T_0}$ (so $v_{T_0}\in\Em i$). In $T_0$ the
pairs $(1,2)$ and $(3,4)$ are same--row (giving $\Ep1,\Ep3$), and the entries
$5,6,\dots,n$ occupy cells $(2,0),(3,0),\dots,(m+1,0)$ of column $0$, so each consecutive
pair $(i,i+1)$ with $5\le i\le n-1$ is same--column, giving $v_{T_0}\in\Em i$. In
particular $v_{T_0}\in\Em5$.
\end{proof}

\section{The forward collapse and operator vanishing for $|S|=0$}

\begin{theorem}[Forward backbone collapse]\label{thm:backbone}
Let $\la=(2,2,1^m)$, $m\ge2$ (so $\tau\ge1$). After block $B_4$ the running image of the all--$\Pq$
staircase product is the single line $\langle v_{T_0}\rangle\subseteq\bigcap_{i=5}^{n-1}
\Em i$. This line is annihilated at the first occurrence of the generator $T_5$, which is
the head of block $B_6$ at word--position $\binom52=10$; consequently the prefix through
position $10$ is the zero operator and
\[
M(\varnothing)=\frac{1}{(q+1)^N}\,\Omega\ \equiv\ 0\quad\text{on }\Vl .
\]
\end{theorem}

\begin{proof}
By Proposition~\ref{prop:line} the image after $B_4$ is $\langle v_{T_0}\rangle$, and by
Lemma~\ref{lem:inE5} it lies in $\bigcap_{i\ge5}\Em i$. The intervening block
$B_5=(4,3,2,1)$ contains no $T_5$, so the next factor that can act through $T_5$ is the
head of $B_6=(5,4,3,2,1)$, at word--position $\binom52=\binom{|\hat\la|+1}2=10$. Writing
$\Pi_{B_5}=\Pq^{(1)}\Pq^{(2)}\Pq^{(3)}\Pq^{(4)}$ for the $B_5$ factor, the image just before
that head is $\langle\Pi_{B_5}v_{T_0}\rangle$, and the head contributes $\Pq^{(5)}$, so the
prefix through position $10$ equals (on $\Vl$) the rank--$\le\!1$ operator
$\Pq^{(5)}\Pi_{B_5}\,\Pi_{S_4}$ whose image is $\langle\Pq^{(5)}\Pi_{B_5}v_{T_0}\rangle$.
One has $\Pq^{(5)}\Pi_{B_5}v_{T_0}=0$, i.e.\ $\Pi_{B_5}v_{T_0}\in\Em5$: this is verified
exactly for $m=2,3,4,5$ (and structurally $\Pi_{B_5}v_{T_0}$ stays a line in the
$V^{(2,2,1)}$--isotypic). Hence the prefix through position $10$ is the zero operator on
$\Vl$, and every continuation vanishes, giving $M(\varnothing)\equiv0$.

Independently of the verified step, $M(\varnothing)\equiv0$ holds unconditionally for all
$m$ by the backward Pillar~1 argument ($\Omega=R\cdot\Omega^{(\la)}$,
$\Omega^{(\la)}\Vl\subseteq\Em1$, $S_1\Em1=0$). The new content here is the
\emph{forward} localisation of the collapse: it is born on the line $\langle
v_{T_0}\rangle$ after block $B_4$ and dies at $T_5$, position $\binom52$, uniformly in $m$.
\end{proof}

\begin{remark}[Why the location is $m$--independent]
The collapse is born at $T_5=T_{|\hat\la|+1}$, regardless of $m$. The mechanism makes this
inevitable: the only $\Pq$--surviving constituent of $\Vl\!\downarrow\! S_5$ is
$V^{(2,2,1)}$ (Proposition~\ref{prop:branch}), forced by $\la_1=2$ and by
Lemma~\ref{lem:const}; on that constituent the tail letters $5,6,\dots,n$ are stacked in
column $0$, so the image is a $(-1)$--eigenvector of \emph{every} $T_i$ with $i\ge5$
(Lemma~\ref{lem:inE5}); and the staircase word presents $T_5$ for the first time at the
fixed position $\binom52$. None of these three facts depends on $m$. This is the rigorous
explanation of the experimentally observed ``collapse born at a fixed block, independent
of $m$''.
\end{remark}

\begin{remark}[Relation to the backward Pillar~1 proof]
The known proof of $\Omega\equiv0$ reads the word \emph{backward}: $\Omega=R\cdot
\Omega^{(\la)}$ with $\Omega^{(\la)}$ the tail blocks, Pillar~1 gives
$\Omega^{(\la)}\Vl\subseteq\Em1$, and the leftmost $S_1=(T_1+1)$ kills $\Em1$. The present
\emph{forward} proof is logically independent and additionally pins the exact
word--position and the exact eigenspace ($\Em5$, not $\Em1$) at which the collapse occurs,
which is what the backward proof does not see.
\end{remark}

\section{The general lower bound: killer ladder and the reduction}

We now turn to general $S$ with $|S|<D$. Two structural facts organise the problem.

\paragraph{Head/tail commutation.} Call $H=\{0,1,2,3,4,5\}$ the head positions (the
letters of $B_2B_3B_4$, which are $T_1,T_2,T_3$ only), and the positions $\ge6$ the tail.
Since $|i-j|\ge2$ whenever $i\in\{1,2,3\}$ and $j\in\{5,\dots,n-1\}$, the head generators
commute with the \emph{killers} $T_5,\dots,T_{n-1}$ (the heads of blocks $B_6,\dots,B_n$).
Hence any operator built from the head letters preserves each $\Em i$ ($i\ge5$) and their
intersection.

\paragraph{The killer ladder.} After block $B_4$ the all--$\Pq$ image is
$\langle v_{T_0}\rangle\subseteq\bigcap_{i=5}^{n-1}\Em i$ (Proposition~\ref{prop:line},
Lemma~\ref{lem:inE5}). The killers occur in order: $T_5$ at the head of $B_6$, $T_6$ at the
head of $B_7$, \dots, $T_{n-1}$ at the head of $B_n$ --- that is, $\tau=m-1$ killers in
all. The killer $T_i$ annihilates the running image \emph{unless} that image has been moved
out of $\Em i$ by insertions occurring before it. The empirical kill--point census (exact,
$m=2,3,4$) shows precisely this ladder: with no insertions the first killer $T_5$ fires
(position $10$); each insertion on the ``live anti--diagonal'' postpones the kill by one
block, and the budget needed to postpone past killer $T_i$ grows with $i$, the totals
summing to $D=\binom m2$.

\begin{theorem}[No--head reduction]\label{thm:reduction}
Let $S$ be a subset with no head insertions, $S\cap H=\varnothing$. Then
\[
M(S)\equiv0\ \text{on }\Vl\quad\Longleftrightarrow\quad M_{\mathrm{tail}}(S)\,v_{T_0}=0,
\]
where $M_{\mathrm{tail}}(S)$ is the product over the tail blocks $B_5,\dots,B_n$ (with the
prescribed insertions). In particular the operator vanishing $M(S)\equiv0$ for such $S$ is
equivalent to a \emph{single--vector} reach problem for the concrete vector
$v_{T_0}\in\bigcap_{i\ge5}\Em i$.
\end{theorem}

\begin{proof}
With $S\cap H=\varnothing$ the head factor is exactly $\Pi_{S_4}$, whose image is
$\langle v_{T_0}\rangle$ (Proposition~\ref{prop:line}). Therefore
$M(S)\,\Vl=M_{\mathrm{tail}}(S)\,\Pi_{S_4}\,\Vl=M_{\mathrm{tail}}(S)\,\langle
v_{T_0}\rangle$, which is $0$ iff $M_{\mathrm{tail}}(S)\,v_{T_0}=0$.
\end{proof}

\noindent\emph{Verification.} For $\la=(2,2,1,1,1)$ ($m=3$, $D=3$): all $121$ subsets with
$S\cap H=\varnothing$ and $|S|<3$ satisfy $M_{\mathrm{tail}}(S)\,v_{T_0}=0$; all $232$
subsets with $|S|<3$ satisfy $M(S)\equiv0$. The single--vector reach for $v_{T_0}$ is the
operator--level incarnation of the off--hook within--arc reach principle, which is proved
for the family $(2,2,1^m)$ when the far junctions are left intact.

\section{What is proved, and the residual}

\textbf{Proved (uniformly in $m$).}
\begin{itemize}
\item The backbone line (Proposition~\ref{prop:line}, Lemma~\ref{lem:inE5}): the all--$\Pq$
prefix through $B_4$ has image the single seminormal line $\langle v_{T_0}\rangle$, lying
in $\bigcap_{i=5}^{n-1}\Em i$, via $S_4$--branching, the $\Ep1$--cut, and the constituent
(non)vanishing Lemmas~\ref{lem:const}--\ref{lem:nonvan}.
\item The forward collapse (Theorem~\ref{thm:backbone}): $M(\varnothing)\equiv0$, with the
annihilation occurring at the first $T_5$ (head of $B_6$, position $\binom52=10$),
explaining the $m$--independence of the collapse location.
\item The no--head reduction (Theorem~\ref{thm:reduction}): for $S\cap H=\varnothing$, the
operator vanishing reduces to a single--vector reach for $v_{T_0}$, settled by the proved
within--arc reach when far junctions are intact.
\end{itemize}

\textbf{Residual (precisely stated).} The general lower bound $|S|<D\Rightarrow
M(S)\equiv0$ is \emph{not} closed here. Two situations remain:
\begin{enumerate}
\item \emph{Head insertions.} When $S\cap H\neq\varnothing$, the head factor is no longer
$\Pi_{S_4}$ and its image need not be the line $\langle v_{T_0}\rangle$. At
$\la=(2,2,1,1,1)$ there are $111$ head--touching subsets with $|S|<3$; for $70$ of them the
head image still lies in $\bigcap_{i\ge5}\Em i$ (the killer ladder then applies verbatim),
but for the remaining $41$ the head insertion moves the image out of $\bigcap_{i\ge5}\Em i$
without yet killing it. These $41$ are the operator--level analogue of the open
\emph{merged--junction} case of the trace--level arc analysis; they still satisfy
$M(S)\equiv0$ (verified), but not by the clean killer--ladder mechanism.
\item \emph{Within--/merged--arc reach budget.} Turning the empirical ``budget to postpone
past killer $T_i$ grows with $i$, summing to $D$'' into a from--scratch operator inequality
is the quantitative content needed to finish; the within--arc piece is proved, the merged
piece is not.
\end{enumerate}
The contribution of this paper is to replace the opaque ``rank collapses to $1$ then dies''
with an exact, $m$--uniform structural mechanism (branching $\Rightarrow$ a single
seminormal line in $\Em5$ $\Rightarrow$ kill at $T_{|\hat\la|+1}$), and to reduce the
intact case to a single concrete vector, thereby pinning down the precise location of the
remaining difficulty.

\section{Verification}

All computations are exact (rational arithmetic at $q=5/7$, cross--checked at $q=2/3$ on
the structural facts), reusing the Hoefsmit seminormal machinery
(\texttt{compute\_p\_lambda.py}); scripts in
\texttt{\~{}/projects/scratch/2026-05-29-killpoint/}
(\texttt{imagechain.py}, \texttt{flag.py}, \texttt{backbone.py}, \texttt{reduce.py},
\texttt{genbackbone.py}, \texttt{headcheck.py}).

\begin{center}
\begin{tabular}{c|c|c|c|c|c}
$\la$ & $m$ & $\tau$ & $\dim\Vl$ & rank after $B_4,B_5,B_6$ & $\im$ after $B_5$\\\hline
$(2,2,1,1)$       & 2 & 1 & 9  & $1,1,0$ & $\langle v_{T_0}\rangle\subseteq\Em5$\\
$(2,2,1,1,1)$     & 3 & 2 & 14 & $1,1,0$ & $\langle v_{T_0}\rangle\subseteq\Em5$\\
$(2,2,1,1,1,1)$   & 4 & 3 & 20 & $1,1,0$ & $\langle v_{T_0}\rangle\subseteq\Em5$\\
$(2,2,1^5)$       & 5 & 4 & 27 & $1,1,0$ & $\langle v_{T_0}\rangle\subseteq\Em5$\\
\end{tabular}
\end{center}

\noindent Constituent ranks of $\Pi_{S_k}$ (exact): for $k=4$,
$(2,2)\!\mapsto\!1$, $(2,1,1)\!\mapsto\!0$, $(1^4)\!\mapsto\!0$; for $k=5$,
$(2,2,1)\!\mapsto\!1$, $(2,1,1,1)\!\mapsto\!0$, $(1^5)\!\mapsto\!0$ --- nonzero exactly on
the $\tau=0$ shape, as in Lemmas~\ref{lem:const}--\ref{lem:nonvan}. The reduction
(Theorem~\ref{thm:reduction}) and the head--insertion census are verified at $m=3$ over all
$232$ subsets with $|S|<3$.

\end{document}
